\documentclass[10pt,a4paper]{amsart}
\usepackage{amsmath,amssymb,amsthm,graphicx}
\usepackage{verbatim}
\newenvironment{reference}{\comment}{\endcomment}
%\newenvironment{reference}{}{}
\newenvironment{slogan}{\comment}{\endcomment}
\newenvironment{history}{\comment}{\endcomment}

% For commutative diagrams we use Xy-pic
\usepackage[all]{xy}

% We use 2cell for 2-commutative diagrams.
\xyoption{2cell}
% Theorem environments.
%
\theoremstyle{plain}
\newtheorem{theorem}[subsection]{Theorem}
\newtheorem{proposition}[subsection]{Proposition}
\newtheorem{lemma}[subsection]{Lemma}

\theoremstyle{definition}
\newtheorem{definition}[subsection]{Definition}
\newtheorem{example}[subsection]{Example}
\newtheorem{exercise}[subsection]{Exercise}
\newtheorem{situation}[subsection]{Situation}

\theoremstyle{remark}
\newtheorem{remark}[subsection]{Remark}
\newtheorem{remarks}[subsection]{Remarks}

\numberwithin{equation}{subsection}

% Macros
%
\def\lim{\mathop{\mathrm{lim}}\nolimits}
\def\colim{\mathop{\mathrm{colim}}\nolimits}
\def\Spec{\mathop{\mathrm{Spec}}}
\def\Hom{\mathop{\mathrm{Hom}}\nolimits}
\def\Ext{\mathop{\mathrm{Ext}}\nolimits}
\def\SheafHom{\mathop{\mathcal{H}\!\mathit{om}}\nolimits}
\def\SheafExt{\mathop{\mathcal{E}\!\mathit{xt}}\nolimits}
\def\Sch{\mathit{Sch}}
\def\Mor{\mathop{\mathrm{Mor}}\nolimits}
\def\Ob{\mathop{\mathrm{Ob}}\nolimits}
\def\Sh{\mathop{\mathit{Sh}}\nolimits}
\def\NL{\mathop{N\!L}\nolimits}
\def\CH{\mathop{\mathrm{CH}}\nolimits}
\def\proetale{{pro\text{-}\acute{e}tale}}
\def\etale{{\acute{e}tale}}
\def\QCoh{\mathit{QCoh}}
\def\Ker{\mathop{\mathrm{Ker}}}
\def\Im{\mathop{\mathrm{Im}}}
\def\Coker{\mathop{\mathrm{Coker}}}
\def\Coim{\mathop{\mathrm{Coim}}}

% Boxtimes
%
\DeclareMathSymbol{\boxtimes}{\mathbin}{AMSa}{"02}

%
% Macros for moduli stacks/spaces
%
\def\QCohstack{\mathcal{QC}\!\mathit{oh}}
\def\Cohstack{\mathcal{C}\!\mathit{oh}}
\def\Spacesstack{\mathcal{S}\!\mathit{paces}}
\def\Quotfunctor{\mathrm{Quot}}
\def\Hilbfunctor{\mathrm{Hilb}}
\def\Curvesstack{\mathcal{C}\!\mathit{urves}}
\def\Polarizedstack{\mathcal{P}\!\mathit{olarized}}
\def\Complexesstack{\mathcal{C}\!\mathit{omplexes}}
% \Pic is the operator that assigns to X its picard group, usage \Pic(X)
% \Picardstack_{X/B} denotes the Picard stack of X over B
% \Picardfunctor_{X/B} denotes the Picard functor of X over B
\def\Pic{\mathop{\mathrm{Pic}}\nolimits}
\def\Picardstack{\mathcal{P}\!\mathit{ic}}
\def\Picardfunctor{\mathrm{Pic}}
\def\Deformationcategory{\mathcal{D}\!\mathit{ef}}
\usepackage{cleveref}
\crefname{originaldefinition}{original definition}{original definitions}
\Crefname{originaldefinition}{Original definition}{Original definitions}
\makeatletter
\newlabel{definition-filtered}{{0121}{source}{}{sourcepointer-1}{}}
\newlabel{definition-filtered@cref}{{[originaldefinition][0][]0121}{[1][0][]source}{}{}{}}
\makeatother
\crefname{originalequation}{original equation}{original equations}
\Crefname{originalequation}{Original equation}{Original equations}
\makeatletter
\newlabel{equation-filtration-cohomology}{{012R}{source}{}{sourcepointer-2}{}}
\newlabel{equation-filtration-cohomology@cref}{{[originalequation][0][]012R}{[1][0][]source}{}{}{}}
\makeatother
\crefname{originallemma}{original lemma}{original lemmas}
\Crefname{originallemma}{Original lemma}{Original lemmas}
\makeatletter
\newlabel{lemma-filtered-complex-ss-converges}{{012V}{source}{}{sourcepointer-3}{}}
\newlabel{lemma-filtered-complex-ss-converges@cref}{{[originallemma][0][]012V}{[1][0][]source}{}{}{}}
\makeatother
\crefname{originalequation}{original equation}{original equations}
\Crefname{originalequation}{Original equation}{Original equations}
\makeatletter
\newlabel{equation-on-top-bigraded}{{012S}{source}{}{sourcepointer-4}{}}
\newlabel{equation-on-top-bigraded@cref}{{[originalequation][0][]012S}{[1][0][]source}{}{}{}}
\makeatother
\crefname{originalequation}{original equation}{original equations}
\Crefname{originalequation}{Original equation}{Original equations}
\makeatletter
\newlabel{equation-at-bottom-bigraded}{{012T}{source}{}{sourcepointer-5}{}}
\newlabel{equation-at-bottom-bigraded@cref}{{[originalequation][0][]012T}{[1][0][]source}{}{}{}}
\makeatother
\crefname{originallemma}{original lemma}{original lemmas}
\Crefname{originallemma}{Original lemma}{Original lemmas}
\makeatletter
\newlabel{lemma-relate-boundedness}{{0BDV}{source}{}{sourcepointer-6}{}}
\newlabel{lemma-relate-boundedness@cref}{{[originallemma][0][]0BDV}{[1][0][]source}{}{}{}}
\makeatother
\crefname{originallemma}{original lemma}{original lemmas}
\Crefname{originallemma}{Original lemma}{Original lemmas}
\makeatletter
\newlabel{lemma-characterize-weak-serre-subcategory}{{0754}{source}{}{sourcepointer-7}{}}
\newlabel{lemma-characterize-weak-serre-subcategory@cref}{{[originallemma][0][]0754}{[1][0][]source}{}{}{}}
\makeatother
\title{Boundedness and convergence of the spectral sequence of a finitely filtered complex}
\author{The Stacks Project Authors}
\date{stacks-project:homology.tex:local-snapshot-20261012; GNU FDL 1.2 or later}
\begin{document}
\maketitle
\noindent\textit{Editorial note.} All delivered mathematical text is the author's own, copied from the pinned Stacks Project chapter file homology.tex; no statement or proof step was rewritten, reordered or removed.\par\medskip
\noindent\textit{Editorial note.} The proof is the author's complete argument for all four assertions; the sentence 'Part (2) is clear from Equation ...' refers to a displayed identity established earlier in the chapter rather than declining a step, and it is preserved verbatim together with the rest of the proof.\par\medskip
\noindent\textit{Editorial note.} Every labelled reference occurring inside the counted spans is carried in assembly.external\_references as a non-load-bearing navigation pointer (same-chapter ones at their definition site, cross-chapter ones at the invocation inside the counted span), so no load-bearing mathematics is replaced by a stub. No bibliography citation occurs inside the counted spans.\par\medskip
\noindent\textit{Editorial note.} The same-name check for this card scanned every staged Stacks chapter of this repository -- all sources/stacks-*/pinned-original-source.tex files, the card's own chapter excluded -- and the labels found are recorded in dedup.same\_name\_cross\_chapter on the card.\par\medskip
\noindent\textit{Editorial note.} The preamble is the chapter's own native preamble excerpt already used for this source: the comment-type reference, slogan and history environments (native preamble lines 8-13), the Xy-pic setup (native preamble lines 14-18) and the native theorem declarations with the mathematical macros (native preamble lines 154-217, the actual last line of the file).\par\medskip
\noindent\textit{Allowed standard prerequisites (editorial).}\begin{itemize}
\item Filtered complexes and the filtration on cohomology. Same chapter, Definition definition-filtered and the construction of the filtration on H\textasciicircum{}n(K). Supplies the notation of the statement and the identity from which finiteness of the filtration on cohomology is read off.
\item The weak convergence criterion. Same chapter, Lemma lemma-filtered-complex-ss-converges together with the displayed descriptions of the top and bottom pieces. Reduces convergence to two displayed equalities, which the proof verifies using the finiteness of the filtration.
\item Boundedness implies regularity and coregularity. Same chapter, Lemma lemma-relate-boundedness. Supplies that a bounded spectral sequence is regular, which gives convergence, and coregular, which lets the author replace a high page by the limiting page.
\item Weak Serre subcategories are closed under extensions. Same chapter, Lemma lemma-characterize-weak-serre-subcategory. Supplies that the middle term of a short exact sequence with outer terms in a weak Serre subcategory again lies in it, applied to the finite filtration of H\textasciicircum{}n.
\end{itemize}\medskip
\ifdefined\hypertarget\hypertarget{sourcepointer-1}{}\fi
\noindent\textit{Source reference.} Original definition 0121: sources/stacks-homology/pinned-original-source.tex, source lines 4884--4884. Navigation pointer invoked at proof line 6582; the counted span contains the pointer only and the cited result is not reproduced inside it, so the pointer is registered as a non-load-bearing navigation stub in the established form for Stacks material. No bibliography citation is involved.\par\medskip
\ifdefined\hypertarget\hypertarget{sourcepointer-2}{}\fi
\noindent\textit{Source reference.} Original equation 012R: sources/stacks-homology/pinned-original-source.tex, source lines 6392--6392. Navigation pointer invoked at proof line 6596; the counted span contains the pointer only and the cited result is not reproduced inside it, so the pointer is registered as a non-load-bearing navigation stub in the established form for Stacks material. No bibliography citation is involved.\par\medskip
\ifdefined\hypertarget\hypertarget{sourcepointer-3}{}\fi
\noindent\textit{Source reference.} Original lemma 012V: sources/stacks-homology/pinned-original-source.tex, source lines 6559--6559. Navigation pointer invoked at proof line 6597; the counted span contains the pointer only and the cited result is not reproduced inside it, so the pointer is registered as a non-load-bearing navigation stub in the established form for Stacks material. No bibliography citation is involved.\par\medskip
\ifdefined\hypertarget\hypertarget{sourcepointer-4}{}\fi
\noindent\textit{Source reference.} Original equation 012S: sources/stacks-homology/pinned-original-source.tex, source lines 6451--6451. Navigation pointer invoked at proof line 6565; the counted span contains the pointer only and the cited result is not reproduced inside it, so the pointer is registered as a non-load-bearing navigation stub in the established form for Stacks material. No bibliography citation is involved.\par\medskip
\ifdefined\hypertarget\hypertarget{sourcepointer-5}{}\fi
\noindent\textit{Source reference.} Original equation 012T: sources/stacks-homology/pinned-original-source.tex, source lines 6461--6461. Navigation pointer invoked at proof line 6565; the counted span contains the pointer only and the cited result is not reproduced inside it, so the pointer is registered as a non-load-bearing navigation stub in the established form for Stacks material. No bibliography citation is involved.\par\medskip
\ifdefined\hypertarget\hypertarget{sourcepointer-6}{}\fi
\noindent\textit{Source reference.} Original lemma 0BDV: sources/stacks-homology/pinned-original-source.tex, source lines 6509--6509. Navigation pointer invoked at proof line 6609; the counted span contains the pointer only and the cited result is not reproduced inside it, so the pointer is registered as a non-load-bearing navigation stub in the established form for Stacks material. No bibliography citation is involved.\par\medskip
\ifdefined\hypertarget\hypertarget{sourcepointer-7}{}\fi
\noindent\textit{Source reference.} Original lemma 0754: sources/stacks-homology/pinned-original-source.tex, source lines 2366--2366. Navigation pointer invoked at proof line 6619; the counted span contains the pointer only and the cited result is not reproduced inside it, so the pointer is registered as a non-load-bearing navigation stub in the established form for Stacks material. No bibliography citation is involved.\par\medskip
\begin{lemma}
\label{lemma-biregular-ss-converges}
Let $\mathcal{A}$ be an abelian category. Let $(K^\bullet, F)$ be a
filtered complex of $\mathcal{A}$. Assume that the filtration on each $K^n$
is finite (see Definition \ref{definition-filtered}). Then
\begin{enumerate}
\item the spectral sequence associated to $(K^\bullet, F)$ is bounded,
\item the filtration on each $H^n(K^\bullet)$ is finite,
\item the spectral sequence associated to $(K^\bullet, F)$ converges
to $H^*(K^\bullet)$,
\item if $\mathcal{C} \subset \mathcal{A}$ is a weak Serre subcategory
and for some $r$ we have $E_r^{p, q} \in \mathcal{C}$ for all
$p, q \in \mathbf{Z}$, then $H^n(K^\bullet)$ is in $\mathcal{C}$.
\end{enumerate}
\end{lemma}

\begin{proof}
Part (1) follows as $E_0^{p, n - p} = \text{gr}^p K^n$.
Part (2) is clear from Equation (\ref{equation-filtration-cohomology}).
We will use Lemma \ref{lemma-filtered-complex-ss-converges} to prove
that the spectral sequence weakly converges. Fix $p, n \in \mathbf{Z}$.
The right hand side of (\ref{equation-on-top-bigraded})
is equal to $F^pK^n \cap \Ker(d) + F^{p + 1}K^n$ because
$F^{p + r}K^n = 0$ for $r \gg 0$. Thus (\ref{equation-on-top-bigraded})
is an equality. The left hand side of (\ref{equation-at-bottom-bigraded})
is equal to $F^pK^n \cap \Im(d) + F^{p + 1}K^n$ because
$F^{p - r + 1}K^{n - 1} = K^{n - 1}$ for $r \gg 0$.
Thus (\ref{equation-at-bottom-bigraded}) is an equality. Since the filtration
on $H^n(K^\bullet)$ is finite by (2) we see that we have abutment.
To prove we have convergence we have to show the spectral sequence is
regular which follows as it is bounded
(Lemma \ref{lemma-relate-boundedness}) and we have
to show that $H^n(K^\bullet) = \lim_p H^n(K^\bullet)/F^pH^n(K^\bullet)$
which follows from the fact that the filtration on $H^*(K^\bullet)$
is finite proved in part (2).

\medskip\noindent
Proof of (4). Assume that for some $r \geq 0$ we have
$E_r^{p, q} \in \mathcal{C}$ for some weak Serre subcategory
$\mathcal{C}$ of $\mathcal{A}$. Then $E_{r + 1}^{p, q}$ is
in $\mathcal{C}$ as well, see
Lemma \ref{lemma-characterize-weak-serre-subcategory}.
By boundedness proved above (which implies that the spectral sequence
is both regular and coregular, see Lemma \ref{lemma-relate-boundedness})
we can find an $r' \geq r$ such that $E_\infty^{p, q} = E_{r'}^{p, q}$
for all $p, q$ with $p + q = n$. Thus $H^n(K^\bullet)$ is an object
of $\mathcal{A}$ which has a finite filtration whose graded pieces
are in $\mathcal{C}$. This implies that $H^n(K^\bullet)$ is in $\mathcal{C}$
by Lemma \ref{lemma-characterize-weak-serre-subcategory}.
\end{proof}

\end{document}
